\chapter{W001: finite-state foundations and complete proofs}
\label{app:finite_foundations}
\HistoricalStatus{New finite-dimensional foundation proofs. These statements
require the hypotheses given here; they do not certify the inherited numerical
services, a bulk prefactor, or a thermodynamic limit.}

This appendix supplies the mathematical details used in the logical narrative.
The Majorana ordering and covariance sign are fixed here rather than inferred
from legacy code. The statements concern finite-dimensional Hilbert spaces.
The perturbative part reconstructs the conditional theorem reviewed in W001-T4;
the remaining identities are proved explicitly for the W001-T5 foundation.

% EDITORIAL QUESTION: Why does physical projection preserve a variational upper bound?
\section{Variational states and exact local projection}
\label{app:projection_proof}
\begin{proposition}
Let $H$ be self-adjoint on a finite Hilbert space, with minimum eigenvalue $E_0$.
Every density operator $\rho$ on that space obeys $\Tr\rho H\geq E_0$.
If $P$ is an orthogonal projector and $\Tr P\rho>0$, then
$P\rho P/\Tr P\rho$ is a density operator on $\Ran P$.
\end{proposition}
\begin{proof}
In an eigenbasis $H|j\rangle=E_j|j\rangle$, the diagonal entries
$p_j=\langle j|\rho|j\rangle$ are nonnegative and sum to one. Thus
$\Tr\rho H-E_0=\sum_jp_j(E_j-E_0)\geq0$. For the projection,
$\langle v|P\rho P|v\rangle=\langle Pv|\rho|Pv\rangle\geq0$,
its trace is $\Tr P\rho$, and its range lies in $\Ran P$.
Normalizing proves the statement. Applying the first result to a Hamiltonian
restricted to an identified physical subspace gives its variational bound.
No bound on that physical problem follows for a state outside that subspace.
\end{proof}

Let $\gamma_1=b^x,\gamma_2=b^y,\gamma_3=b^z,\gamma_4=c$ at one site.
Anticommuting distinct Majoranas gives
$D^\dagger=(-1)^6D=D$ and $D^2=(-1)^6\Id=\Id$.
Each Majorana anticommutes with $D$, so every bilinear commutes with $D$.
For $\widetilde\sigma^\alpha=\ii b^\alpha c$, its square is one and
\begin{align}
 \widetilde\sigma^x\widetilde\sigma^y&=b^xb^y,\\
 D\widetilde\sigma^z&=\ii b^xb^yb^zc\,b^zc=-\ii b^xb^y.
\end{align}
Therefore $\widetilde\sigma^x\widetilde\sigma^y=\ii D\widetilde\sigma^z$;
cyclic permutations give the other products. Even bilinears on disjoint sites
commute because interchanging their four constituent pairs gives a positive
sign. These facts prove the Pauli relations on the $D_i=+1$ joint subspace.

To count its dimension, expand $P_G=2^{-N}\prod_i(\Id+D_i)$. Every nonempty
product of $D_i$ has zero trace: conjugating by a Majorana at one included site
changes its sign. Hence $\Tr P_G=2^{-N}4^N=2^N$.
With $f_1=(b^x+\ii b^y)/2$, $f_2=(b^z+\ii c)/2$, one has
$\ii b^xb^y=2n_1-1$ and $\ii b^zc=2n_2-1$. Consequently
$D=-(2n_1-1)(2n_2-1)$ and the physical local subspace is spanned by
$|01\rangle,|10\rangle$ in this ordering.

The projected state of $\rho_{\rm aux}=\Id_4/4$ is
$\rho_G=(|01\rangle\langle01|+|10\rangle\langle10|)/2$.
Both occupations equal one half, all hopping/pairing bilinears have zero mean,
and therefore all entries of $\Omega$ vanish. But $\langle D\rangle=1$.
A Gaussian state with that covariance has vanishing four-distinct-Majorana
expectation, proving that this projected state is not Gaussian.
Finally, for an involution $D$,
$\Tr\rho(\Id-D)=0$ implies no weight in its negative sector, since
$\Id-D=2P_-\succeq0$. This proves the special support statement used in the
main chapter, without asserting it for arbitrary zero-mean constraints.

% EDITORIAL QUESTION: How do positivity, Wick factorization, and purity follow in a canonical mode basis?
\section{Gaussian covariances and the pure-state orbit}
\label{app:gaussian_proof}
The CAR imply that $\Omega$ in Eq.~\eqref{eq:narr:covariance} is real skew.
For any complex vector $z$, positivity of
$\langle(\sum_a z_a\gamma_a)^\dagger(\sum_bz_b\gamma_b)\rangle$
means $\Id-\ii\Omega\succeq0$. A real skew matrix admits an orthogonal canonical
form $O\bigoplus_j\lambda_jJ_2\,O^T$; the eigenvalues of
$\Id-\ii\lambda_jJ_2$ are $1\pm\lambda_j$. Thus any state has
$|\lambda_j|\leq1$.

Conversely, in this basis let independent complex-mode occupations have
probabilities $p_j=(1+\lambda_j)/2$. The product density operator
\begin{equation}
 \rho=\bigotimes_j\bigl((1-p_j)|0_j\rangle\langle0_j|
                         +p_j|1_j\rangle\langle1_j|\bigr)
\end{equation}
is a Gaussian state with that covariance. For $0<p_j<1$ it is the exponential
of a quadratic number operator; endpoints are its pure or partially pure
limits. Orthogonal changes of a finite Majorana basis are unitarily implementable
on the full Fock space (the Pin representation), so the construction applies
to the original basis. It is pure exactly when all $p_j$ are endpoints, which
is equivalent to $\Omega^2=-\Id$.

For completeness the Wick signs can be checked without assuming a spin
factorization. Introduce anticommuting auxiliary variables $\theta_a$. In the
canonical product state, the generating polynomial for ordered distinct
correlations factors into the independent pairs. Only the empty term and the
pair term survive in each mode, so it equals the Grassmann Gaussian with
pair coefficients $C_{ab}=\langle\gamma_a\gamma_b\rangle=-\ii\Omega_{ab}$ for
$a<b$. Expanding this exponential, a coefficient with $2m$ distinct indices
is the signed sum over all pairings,
\begin{equation}
 \Pf C_I=\frac{1}{2^m m!}\sum_{\pi\in S_{2m}}\sgn(\pi)
             \prod_{j=1}^m C_{a_{\pi(2j-1)},a_{\pi(2j)}}.
\end{equation}
The pair contractions transform multilinearly under an orthogonal basis
change, so the same identity holds in every real Majorana basis. It gives
$\Pf(-\ii\Omega_I)=(-\ii)^m\Pf\Omega_I$ as stated. Endpoint limits preserve
this polynomial identity. Multiplying the ordered $2n$-point result by
$(-\ii)^n$ gives $\langle\mathcal P\rangle=(-1)^n\Pf\Omega$.

All pure covariances are orthogonally conjugate to a fixed complex structure.
The orthogonal transformations preserving it are precisely the complex-linear
isometries $U(n)$. The orbit is therefore $O(2n)/U(n)$, of dimension
$n(2n-1)-n^2=n(n-1)$. Since $U(n)$ embeds in $SO(2n)$ and is connected, the
orbit has the two components inherited from $O(2n)$. Squaring
$\Pi=(\Id+\ii\Omega)/2$ using $\Omega^2=-\Id$ gives $\Pi^2=\Pi$; it is
Hermitian and has trace $n$, proving the one-particle projector claim.

% EDITORIAL QUESTION: Why is a covariance commutator a stationarity condition rather than a stability theorem?
\section{Tangent variations and constrained stationarity}
\label{app:stationarity_proof}
Differentiating $\Omega^T=-\Omega$ and $\Omega^2=-\Id$ shows that a tangent
$Y$ satisfies $Y^T=-Y$ and $Y\Omega+\Omega Y=0$. Conversely set
$X=-Y\Omega/2$. Anticommutation gives $X^T=-X$ and
$[X,\Omega]=Y$, so every such tangent arises from the orthogonal orbit.
Using Eq.~\eqref{eq:narr:gradient} and cyclicity,
\begin{equation}
 4\delta\mathcal L=\Tr\bigl((\Omega A_{\mathcal L}^T
                  -A_{\mathcal L}^T\Omega)X\bigr)
             =\Tr([A_{\mathcal L},\Omega]X).
\end{equation}
The commutator of two real skew matrices is real skew. If this trace is zero
for every skew $X$, choosing $X=[A_{\mathcal L},\Omega]^T$ proves the
commutator vanishes. Conversely its vanishing makes every orbit variation
zero. Multipliers require constraint regularity as stated in the main chapter;
the calculation provides no sign for the second variation.

For moment positivity, choose arbitrary coefficients $z_a$ and
$B=\sum_a z_aO_a$. Then $z^\dagger Mz=\Tr\rho B^\dagger B\geq0$.
If $\mathcal M_{\rm phys}\subseteq\mathcal M_{\rm rel}$ and the objective
$\ell$ equals energy on physical moments, then
\begin{equation}
 \inf_{M\in\mathcal M_{\rm rel}}\ell(M)
 \leq\inf_{M\in\mathcal M_{\rm phys}}\ell(M)=E_0.
\end{equation}
A single feasible relaxed point satisfies the opposite-direction inequality
relative to the relaxed infimum, explaining why it alone is not a certified
physical lower bound.

% EDITORIAL QUESTION: What exact assumptions make a Schur reduction and its expansion valid?
\section{Finite Feshbach map, norm control, and second order}
\label{app:feshbach_proof}
% CLAIM: W001-FND-01, W001-FND-02 (note: doc/audits/w001_foundation_review.md).
Let $H_0,V$ be self-adjoint on a finite Hilbert space, $P=P^\dagger=P^2$,
$Q=\Id-P$, $[P,H_0]=0$, and $\lambda$ real. Write an eigenvector of
$H=H_0+\lambda V$ as $\psi=p+q$, $p=P\psi$, $q=Q\psi$. Its block equations are
\begin{equation}
 PHPp+PHQq=Ep,\qquad QHPp+QHQq=Eq.
\end{equation}
If $E-QHQ$ is invertible on $\Ran Q$, the second equation gives
$q=(E-QHQ)^{-1}QHPp$. Substitution gives $F(E,\lambda)p=Ep$ with
Eq.~\eqref{eq:narr:feshbach}. Conversely a nonzero solution $p$ of that
reduced equation reconstructs an eigenvector by this formula. A full eigenvector
in this resolvent domain cannot have $p=0$, since then invertibility forces
$q=0$. No inverse on the ambient $P$ block is being taken.

At $\lambda=0$ set $B=E-QH_0Q$ on $\Ran Q$ and $R_0=B^{-1}$.
Factor, preserving operator order,
\begin{equation}
 B-\lambda QVQ=(\Id-\lambda QVQ R_0)B.
\end{equation}
If $|\lambda|\|QVQ R_0\|<1$, the inverse is
$R_0\sum_{r\geq0}(\lambda QVQ R_0)^r$, convergent in operator norm.
Using $PH_0Q=0$ proves Eq.~\eqref{eq:narr:neumann}. A remainder after $r=s$
has norm bounded by
\begin{equation}
 |\lambda|^2\|PVQ\|\,\|R_0\|\,\|QVP\|
 \frac{(|\lambda|\|QVQ R_0\|)^{s+1}}
      {1-|\lambda|\|QVQ R_0\|}.
\end{equation}
This is a conditional bound for the \emph{fixed-energy reduced operator},
not automatically a remainder for an energy-independent effective Hamiltonian
or for the solution $E(\lambda)$ of the nonlinear spectral equation.

The spectral theorem gives
$\|R_0\|=1/\operatorname{dist}(E,\spec(QH_0Q|_{\Ran Q}))$ for real $E$
outside that spectrum. If its true lower edge is $b$ and $E<b$, this is
$1/(b-E)$. A flux external window below $E_0+\Lambda$ therefore has the bound
$1/(\Delta_{\rm flux}-\Lambda)$ only when the \emph{true chosen physical}
eliminated edge is $E_0+\Delta_{\rm flux}$ and $\Lambda<\Delta_{\rm flux}$.
Replacing this edge by one sampled sector has no such guarantee.

Finally take $P$ to be the complete ground eigenspace of $H_0$ at $E_0$,
with $QH_0Q-E_0>0$ on $\Ran Q$. Then
$R_0=-(QH_0Q-E_0)^{-1}$ and for $\psi\in\Ran P$,
\begin{equation}
 \langle\psi|PVQ R_0 QVP|\psi\rangle
 =-\langle QV\psi|(QH_0Q-E_0)^{-1}|QV\psi\rangle.
\end{equation}
Taking a positive square root proves Eq.~\eqref{eq:narr:second_order} and shows
that the second-order operator vanishes exactly when $QVP=0$. The explicit
two-state example in the main chapter disproves the inference from $PVP=0$.
This sign theorem does not assume a flux projector is the ground projector.

% EDITORIAL QUESTION: Which thermal statements depend on holding the Hamiltonian fixed?
\section{Gibbs identities and the free-energy principle}
\label{app:gibbs_proof}
For finite fixed $H$, differentiating its spectral sum gives
$\partial_\beta\ln Z=-E$ and
$\partial_\beta E=-\langle H^2\rangle+E^2$. Since
$\partial_T=-\beta^2\partial_\beta$, Eq.~\eqref{eq:narr:gibbs} follows.
The identity $\ln\rho_\beta=-\beta H-\ln Z$ gives its entropy formula.
As $\beta\to0$, $\rho_\beta\to\Id/d$ and $S\to\ln d$; use $d=2^N$ for
physical spins and $d=4^N$ for the unprojected four-Majorana space.
A positive finite sum $Z(T)$ is analytic for real $T>0$, so its logarithm is
analytic there. A finite sample consequently has no thermodynamic singularity
at positive real temperature.

For any density $\rho$, the trial free energy satisfies
\begin{equation}
 \Tr\rho H-T S(\rho)-F_\beta
    =T D(\rho\Vert\rho_\beta)\geq0.
\end{equation}
Here is a finite proof of the inequality. Diagonalize $\rho$ and the positive
$\sigma=\rho_\beta$ with eigenvalues $p_i\geq0$, $q_j>0$ and set
$t_{ij}=|\langle i|j\rangle|^2$. Row and column sums of $t$ are one, so
\begin{equation}
 D(\rho\Vert\sigma)=\sum_{ij}t_{ij}
    \left[p_i\ln\frac{p_i}{q_j}-p_i+q_j\right]\geq0.
\end{equation}
The scalar bracket equals $q_j(x\ln x-x+1)$ at $x=p_i/q_j$ and is nonnegative
by convexity, including $x=0$ by continuity. This proves the claim and shows
where the state space enters a thermal variational construction.

For a quadratic block, $Z_j=2\cosh(\beta\varepsilon_j/2)$ and
$p_j=(1+e^{\beta\varepsilon_j})^{-1}$. Thus
$\Omega_j=-\tanh(\beta\varepsilon_j/2)J_2$. Since $(\ii J_2)^2=\Id$,
$\ii\tanh(\ii\beta\varepsilon_jJ_2/2)$ has exactly that value.
Orthogonal covariance proves Eq.~\eqref{eq:narr:thermal_covariance} for all
real skew $A$, including its zero blocks. This proves the fixed-$A$ identity,
not the derivative of a temperature-dependent mean-field functional.
For the binary flux toy, dividing the difference of its two Boltzmann weights
by their sum immediately gives Eq.~\eqref{eq:narr:binary_flux}.

% EDITORIAL QUESTION: How are detailed balance and the integrated spectral weight derived?
\section{Lehmann expansion and exact spectral sum}
\label{app:lehmann_proof}
In an energy eigenbasis,
\begin{equation}
 \Tr\rho_\beta A(t)A=Z^{-1}\sum_{mn}e^{-\beta E_m}
             e^{\ii(E_m-E_n)t}A_{mn}A_{nm}.
\end{equation}
Hermiticity makes $A_{mn}A_{nm}=|A_{mn}|^2$. Fourier transforming the
exponential with the stated convention yields Eq.~\eqref{eq:narr:lehmann}.
Integrating over $\omega/(2\pi)$ and using completeness gives
$Z^{-1}\sum_m e^{-\beta E_m}\langle m|A^2|m\rangle$.
For detailed balance, exchange $m,n$ in the expression for $S(-\omega)$.
On the support $\omega=E_n-E_m$, one has
$e^{-\beta E_n}=e^{-\beta\omega}e^{-\beta E_m}$, proving the distributional
identity in Eq.~\eqref{eq:narr:kms}, including degeneracies and elastic weight.
No statement about a tunneling current follows without a probe model.
